%%%%%%%%%%  scan idx 83  |  folio 70  %%%%%%%%%%
\(z\in\underline K(X)\).
Alors, la formule de Künneth donne \(f_!(x)=\chi(y)z\), où pour un faisceau algébrique cohérent \(F\) sur \(P\), l'on pose~:
\[
  \chi(F)=\sum_{i\geq 0}(-1)^i\dim H^i(P,F).
\]
On est alors aussitôt ramené à prouver la formule
\begin{equation}\tag{2.40}
  \chi(F)=\varkappa_r\bigl(\ch_P(F)\,\mathcal T(P)\bigr)
\end{equation}
où \(\varkappa_r\) désigne la composante de degré \(r=\dim P\). En vertu de la proposition 2.13, on peut supposer que \(F=\mathcal O_Q\), où \(Q\) est une sous-variété linéaire de \(P\). Alors \(\chi(F)=1\), puisque c'est le genre arithmétique de \(Q\)~; reste à voir que le deuxième membre de la formule (2.40) vaut \(1\).

Si en effet \(V\) est un hyperplan de \(P\) et si l'on pose \(\xi=\gamma_P(\mathcal O_V)\) et \(L=\ell(V)\), on aura~:
\[
  1-\xi=\gamma_P(E(-L))
\]
(où \(E(-L)\) est le fibré vectoriel défini par \(-L\)), d'où~:
\[
  C_P(1-\xi)=[1,1-L],
\]
et par suite
\[
  \ch_P(1-\xi)=e^{-L},\qquad\text{d'où}\qquad \ch_P(\xi)=1-e^{-L},
\]
et donc
\[
  \ch_P(\xi^p)=(1-e^{-L})^p.
\]
Or si \(p\) est la dimension de \(Q\) dans \(P\), on a \(\xi^p=\gamma_P(Q)\); on en déduit
\[
  \ch_P(F)\,\mathcal T(P)
  =
  \bigl(L/(1-e^{-L})\bigr)^{r+1}(1-e^{-L})^p
\]
et il faut vérifier que le coefficient de \(L^r\) dans cette série formelle en \(L\) vaut \(1\), ce qui résulte de la propriété caractéristique de la série formelle \(L/(1-e^{-L})\), à savoir que le coefficient de \(L^q\) dans \(\bigl(L/(1-e^{-L})\bigr)^{q+1}\) vaut \(1\) (poser \(q=r+1-p\)).

